% AUTO-GENERATED by scripts/generate_paper_results.py from
% evaluation/claim_validity/v1/results.json. Do not hand-edit.
\begin{table}[t]
\centering
\small
\begin{tabular}{lrr}
\toprule
\textbf{Metric} & \multicolumn{2}{r}{\textbf{Value}} \\
\midrule
DECLARATIVE\_CLAIM precision & \multicolumn{2}{r}{94.9\%} \\
DECLARATIVE\_CLAIM recall & \multicolumn{2}{r}{93.0\%} \\
DECLARATIVE\_CLAIM F1 & \multicolumn{2}{r}{93.9\%} \\
Overall accuracy (N=400) & \multicolumn{2}{r}{94.0\%} \\
Macro-F1 & \multicolumn{2}{r}{94.0\%} \\
Train-template accuracy (N=318) & \multicolumn{2}{r}{94.0\%} \\
Template-held-out accuracy (N=82) & \multicolumn{2}{r}{93.9\%} \\
\bottomrule
\end{tabular}
\hfill
\begin{tabular}{lrrr}
\toprule
\textbf{NON\_CLAIM subtype} & \textbf{N} & \textbf{Rejected} & \textbf{Leaked} \\
\midrule
\texttt{code} & 22 & 22 & 0 \\
\texttt{fragment} & 22 & 20 & 2 \\
\texttt{heading} & 23 & 22 & 1 \\
\texttt{list\_item} & 22 & 22 & 0 \\
\texttt{metadata} & 23 & 23 & 0 \\
\texttt{procedural\_instruction} & 22 & 15 & 7 \\
\texttt{question} & 22 & 22 & 0 \\
\texttt{quote} & 22 & 22 & 0 \\
\texttt{table\_cell} & 22 & 22 & 0 \\
\bottomrule
\end{tabular}
\caption{SMRITI-ClaimValidity-v1 (audit round P1-7, diversity/
template-held-out extension P1-F): Phase 4's real, unmodified
\texttt{AssertionClassifier} scored against 400 construction-defined
DECLARATIVE\_CLAIM/NON\_CLAIM spans (an additional 100 AMBIGUOUS spans
are reported separately as a resolution distribution, never scored
right/wrong). ``Leaked'' counts a genuine non-claim the classifier
nonetheless emitted as \texttt{DECLARATIVE\_ASSERTION}, the
consequential error direction (claim-graph noise), not the reverse.
Train-template/template-held-out accuracy is a diversity/robustness
check, not an ML generalization claim: \texttt{AssertionClassifier}
is a fixed rule-based gate, never trained on this or any other
benchmark.}
\label{tab:claim-validity}
\end{table}


Phase 4's own validator explicitly checks structural well-formedness
(no headings, no metadata, no code reach claim construction) but, by its
own documentation, does not evaluate claim \emph{semantics} --- it can
prove ``I don't emit headings, metadata, code, etc.'' but not ``the
objects I do emit are genuinely valid declarative claims'' (P1-7, audit
round). We close this with SMRITI-ClaimValidity-v1
(\texttt{evaluation/claim\_validity/v1/}): 500 candidate spans with a
construction-defined ground truth, independent of whatever the
classifier itself would decide, scored against Phase 4's real,
unmodified \texttt{AssertionClassifier} (no human annotation, per this
project's established convention for benchmarks with no annotator
available). 200 spans are unambiguous DECLARATIVE\_CLAIM, 200 are
unambiguous NON\_CLAIM spread across all nine non-assertion
\texttt{AssertionType} subtypes the classifier can emit, and 100 are
genuinely AMBIGUOUS (bare gerund phrases like ``Preparing the Dough,''
elliptical clauses like ``Still under investigation'') --- these last
100 are reported only as a resolution distribution, never scored
right/wrong, for the same reason SMRITI-Controlled-v2's
ambiguous\_abstention category is (\S\ref{sec:controlled} finding 9):
there is no single correct label for a genuinely ambiguous span by
construction.

\textbf{Diversity rebuild (P1-F, audit round).} The version of this
benchmark evaluated below is a rebuild: the previous version drew each
NON\_CLAIM subtype's ~22 instances from a fixed pool of only ~10
templates via modulo indexing, so most instances within a category were
\emph{verbatim duplicates} of one of a handful of strings --- a
classifier could pass by memorizing template wording rather than
semantic function. Every category now crosses a substantially larger
template pool with entity substitution wherever the template allows it,
mixes in entity-free generic phrasings, and records which literal
template produced each instance so a template-held-out split (never the
same literal template in both \texttt{train} and
\texttt{template\_heldout}) is possible. Six of the nine NON\_CLAIM
subtypes (heading, list\_item, metadata, procedural\_instruction,
question, table\_cell) now have zero verbatim duplicate instances
(\texttt{unique\_text\_count} = $n$); the remaining three (code,
fragment, quote) still repeat a handful of their larger-but-still-finite
template pools. This changes the reported numbers below from an earlier
draft of this section: a template pool artificially narrow enough to
hide a failure mode by only exercising it once or twice is a benchmark
weakness, not a result worth preserving, so we report the rebuilt
benchmark's real numbers rather than the earlier, less diverse ones.

\textbf{(1) Overall accuracy is 94.0\% (macro-F1 94.0\%) on the 400
scoreable spans (TP=186, FP=10, FN=14, TN=190) --- the classifier
correctly separates non-claim idioms from genuine assertions the large
majority of the time, and the diversity rebuild's lower accuracy than an
earlier, narrower-template draft (96.0\%) reflects a benchmark that now
exercises more of the classifier's decision surface, not a regression in
the classifier itself.} Six of the nine NON\_CLAIM subtypes (code,
list\_item, metadata, question, quote, table\_cell) are rejected
perfectly, 100\% recall each.

\textbf{(2) The dominant leak is the same reproducible parser/classifier
interaction the narrower benchmark draft already found, now surfaced at
higher volume because entity substitution multiplied its instance
count.} 7 of 22 \texttt{procedural\_instruction} spans leak into
DECLARATIVE\_ASSERTION, all conditional imperatives of the same two
template families (``Increase \{entity\}'s timeout if requests keep
failing.,'' ``Disable \{entity\}'s feature flag if issues persist.'')
--- the subordinate \emph{if}-clause's own subject (``requests,''
``issues'') satisfies the classifier's \texttt{has\_subject\_anywhere}
fallback check, which exists to rescue genuine assertions from tokenizer
quirks (\S\ref{sec:system}) but here lets a conditional imperative's
embedded clause masquerade as the main clause having a subject. 2 of 22
\texttt{fragment} spans (``see section see section'') and 1 of 23
\texttt{heading} spans (``The Aerotrix scheduler Release Notes'') leak
the same way, both newly surfaced by the wider template pool rather than
present in the earlier draft.

\textbf{(3) The recall-loss side is broader than a single lexical
ambiguity: 14 DECLARATIVE\_CLAIM instances, spanning five distinct
template families, are misclassified NON\_CLAIM (13 as
\texttt{LIST\_ITEM}, 1 as \texttt{FRAGMENT}), all via the same
\texttt{no\_verb\_root\_bare\_noun\_phrase}/
\texttt{verb\_root\_no\_subject\_not\_imperative} path.} The earlier,
narrower draft found only spaCy's noun/verb tagging of \emph{processes}
(a genuine part-of-speech ambiguity: \emph{processes} as the plural noun
vs. the intended third-person-singular verb); that specific case remains
(2 instances here). The wider template pool now also surfaces
\emph{handles} (4 instances), \emph{runs as a single replica in \ldots}
(2 instances), and two longer-subject templates (``The team responsible
for \{entity\} reported \ldots,'' ``Because it lacks a retry policy,
\{entity\} drops requests \ldots'') failing the same way (1 instance
each) --- the same class of spaCy small-model verb-root-detection
failure as (2), now shown to affect several distinct lexical items and
subject constructions rather than one isolated word.

\textbf{(4) The template-held-out check shows no meaningful gap: 94.0\%
accuracy on templates used in \texttt{train} (N=318) vs.\ 93.9\% on
\texttt{template\_heldout} templates (N=82).} Because
\texttt{AssertionClassifier} is a fixed rule-based gate never trained on
this or any benchmark, this is not an ML generalization result --- there
is no model to overfit --- but it is evidence that (1)--(3)'s failure
modes are genuine linguistic-construction weaknesses of the parser/
classifier pipeline, not artifacts concentrated in whichever templates
happened to be sampled into this particular 400-span draw.

\textbf{(5) On genuinely ambiguous spans, the classifier's bias remains
uniform and one-directional: 99 of 100 resolve to NON\_CLAIM, 1 to
DECLARATIVE\_ASSERTION.} This is a disclosed property of this specific
classifier design, not a claim that NON\_CLAIM is the ``right'' answer
for these spans --- by construction, there is no right answer --- but it
is informative: whatever a future extractor built on this substrate
does with genuinely hedged or elliptical text, this classifier's current
rule set almost never mistakes it for a confident assertion, which is
the safer of the two possible biases for a system whose entire
downstream graph is built from what it lets through.

We read this benchmark as a real, if narrow, semantic validation this
paper did not previously have: 94.0\% accuracy on 400 construction-defined
spans drawn from a substantially more diverse template pool than the
earlier draft, with every failure mode traced to specific, reproducible
causes rather than left as an unexplained residual, template-held-out
evidence that those failure modes are not sampling artifacts, and the
classifier's behavior on genuinely hard cases disclosed as a design
property rather than measured against a label that does not exist.
